
\section{Neural network}

(TODO: rewrite)

Neural network is a widely used machine learning model.
Its popularity is heavily based on its exceptional learning capacity allowing it to learn complex relations between features.
Furthermore, it can and is being employed in almost every subfield of machine learning, including natural language processing, computer vision, robotics, and many more due to the ease with which its architecture can be customized for given task.

\subsection{Definition}

As a neural network consists of multiple layers, it is first necessary to provide a formal definition of these layers.
Due to the various types of these layers (see \ref{ss:LayerTypes} Layer types), the definition is very general.

(TODO: is $\mathbb{R}^0$ ok? layer parts should be differentiable)
\begin{definition}[Neural network layer]
\label{d:layer}
    Neural network layer is a real vector function $f: (\Theta; \mathbb{R}^n) \to \mathbb{R}^m$, where $\Theta \subseteq \mathbb{R}^l$ is the space of learnable parameters, $l \ge 0$ and $n, m \ge 1$.
     
     Then for the learnable parameters $\theta \in \Theta$ and the input $x \in \mathbb{R}^n$, the output of the layer is defined as $f(\theta; x)$. 
\end{definition}

It is important to note that not every type of layer contains learnable parameters.
Finally, the neural network can be defined as the composition of those layers.

\begin{definition}[Neural network]
    Let $f_1, f_2, \dots f_n$, $f_i : (\Theta_i, \mathbb{R}^{m_i}) \to \mathbb{R}^{m_{i + 1}}$, $m_1, \dots, m_{n+1} \in \mathbb{N}$ be $n \ge 1$ neural network layers.
    Then we define the neural network as their composition $\mathbf{f}$, formally:
    $$\mathbf{f} \coloneqq f_n \circ \dots  \circ  f_2 \circ f_1.$$
    $f_1$ is the input layer, $f_n$ is the output layer and all layers in between them are hidden layers. 
\end{definition}


\subsection{Layer types}
\label{ss:LayerTypes}
(TODO: some of those layers do not have to be defined formally, it would maybe help the readability)

\subsubsection{Linear layer}

\begin{definition}[Linear layer]
 Suppose that we have functions $f_1, f_2, \cdots, f_n$, where $f_i : \mathbb{R}^{m} \times \mathbb{R} \times \mathbb{R}^{m} \to \mathbb{R}$, for $i \in \{1, 2, \dots, n\}$, furthermore, these functions are in the form $f_i(w, b; x) \coloneqq w^Tx + b$, where $w \in \mathbb{R}^m$ are the learnable weights, $b \in \mathbb{R}$ is the (learnable) intercept, and $x \in \mathbb{R}^m$ is the input vector.
    Then the linear layer (sometimes referred to as dense) $\mathbf{f}$ is the multiple of these functions $\mathbf{f} \coloneqq (f_1, f_2, \dots f_n)$.
\end{definition}

\subsubsection{Activation functions}

Activation functions are rarely considered as layers; however, they behave similarly to other types of layers, and the general definition \ref{d:layer} allows us to define them as one.

\begin{definition}[Traditional activation functions]
\label{d:activation}
 The activation layer is a real vector function $\mathbf{f}: \mathbb{R}^m \to \mathbb{R}^m$ in a form of $\mathbf{f}(x) = \left(f(x_1), f(x_2), \dots, f(x_m)\right)$, where $x = (x_1, x_2, \dots, x_m) \in \mathbb{R}^m$ is the input and $f:\mathbb{R}\to\mathbb{R}$ is a simple function (TODO: simple is perhaps not the word).
\end{definition}

The traditional choices for $f$ (from definition \ref{d:activation}) in are the following:
$$\text{ReLU}(x) \coloneqq \max(0, x);$$
$$\text{SELU} \coloneqq \lambda \begin{cases}
    \alpha \left( e^x - 1\right) & x < 0 \\ x & x \ge 0
\end{cases}\ , \text{where}\ \lambda \approx 1.0507\ \text{and}\ \alpha \approx 1.67326 \ .$$

\subsubsection{2D convolution}
(TODO: rewrite)
The operation of convolution has its roots in physical signal processing (TODO: source).
Its main goal is to smooth out the noisy signal; how the smoothing is exactly done is controlled by the kernel choice.
Convolution is usually denoted by an asterisk and its formal definition is as follows:
$$(x \star w)(t) = \int x(a)s(t-a) \mathrm{d}a \ .$$
This relation can be straightforwardly converted to multidimensional case.
Then its multidimensional discrete form is:
$$(I \star K) = \sum_m\sum_nI(m, n)K(i-m, j-n) \ .$$
However, this version is not implemented in machine learning libraries, e.\ g.\ PyTorch.
Instead, as convolution is meant, cross-correlation defined as:
$$(K \star I) = \sum_m\sum_nI(i+m, j+n)K(m, n) \ .$$
This change is equivalent to flipping the kernel. \cite{Goodfellow2016}

\subsubsection{Pooling}

\subsection{Loss function}

The key component of every possible neural network is its loss function.
Determines the whole learning process, quantifies the model performance, and shows its biggest weaknesses, where the model needs to be improved.
Then the definition of a training process is equivalent to minimizing the loss function.
There are countless possible functions; however, the mean squared error (MSE), Mean Absolute Error (MAE), and Root MSE (RMSE) are the most common ones. (TODO: maybe include PSNR and SSIM)
Their value depends on the true value $Y$ and the model output $\hat{Y}$.
$$
\text{MSE}(Y, \hat{Y}) \coloneqq \frac{1}{n} \sum_{i=1}^n \left( \hat{Y}_i - Y_i \right)^2
$$
$$
\text{MAE}(Y, \hat{Y}) \coloneqq \frac{1}{n} \sum_{i=1}^n  \left| \hat{Y}_i - Y_i \right|$$
$$
\text{RMSE}(Y, \hat{Y}) \coloneqq \sqrt{\text{MSE}(Y, \hat{Y})}
$$

As stated above, it is possible to modify these functions by completely replacing them or simply adding one or more loss terms.
The latter approach will be described in more detail in \ref{s:physical_constraints}.
The input of the model may also be incorporated into the loss function. (TODO: ref somewhere)

\subsection{Optimizer}

The optimizer specifies the learning strategy. 
It decides how the model will react and adjust its learnable parameters to performance changes, measured by loss function.
(TODO: too many 'used' in the whole thesis, find better synonyms)
In early stages of neural networks, usually stochastic gradient descent was used.
Later, an approach inspired by ball movement on a mountainous land (graph of loss function), known as stochastic gradient decent with momentum followed.
Nevertheless, both of them struggle to get out of local minima.
Recently, the algorithm known as Adam dominates this field. (TODO: source, maybe also pseudocode?)

%===================================================================================================

Single image super-resolution (SISR) refers to the problem of transforming a low resolution (LR) input image into its high resolution (HR) version.
Obviously, there are many possibilities for such a function; hence it implements a mapping from low-dimensional to high-dimensional space.
Applying one of these functions to a LR image yields a reconstructed image, known as a super-resolution (SR) image.
The term "downscaling" denotes this procedure.
In the early stage of research around this topic, mainly statistical methods were used to achieve the desired mapping. 
Due to advances in hardware and machine learning methods, deep learning (DL) models were employed. (TODO: Harder et al. have good SISR introduction)

% The incorporation of physical constraints into machine learning models is commonly categorized into two principal approaches: soft and hard physical constraints.  
% Soft physical constraints are implemented by augmenting the loss function such that the model is penalized for violations of the underlying physical laws. In this framework, the optimization procedure seeks to minimize not only the data-driven prediction error but also a term quantifying the degree of physical inconsistency.  
% In contrast, hard physical constraints are enforced by designing the model architecture so that any admissible output is, by construction, consistent with the prescribed physical principles. In this case, the structure of the model precludes the generation of physically infeasible solutions. (TODO: source)
% ~\cite{Harder2023}


In the following soft constraints, the term \textit{dilated SR derivative} will be used.
It is expected that, through downscaling, the derivatives of input LR images will be preserved in output SR images.
Those derivatives in SR images are calculated in a greater distance than one, as in LR neighboring pixels.
Moreover, for scale factor $s$, there are $s^2$ corresponding SR differences to one LR difference. (FIXME: very bad formulation)
Formally, assume we have pixels indexed as $I^\text{LR}_{i,j}$ for LR pixels in the $i$-th row and $j$-th column, and SR pixels downscaled from $I^\text{LR}_{i,j}$ indexed as 
$\left\{ I^\text{SR}_{i,j,k} \mid k \in \left\{1, 2, \dots \text{scale\_factor}^2 \right\}  \right\}$.
Then the mean SR derivative in $i$-th row and $j$-th column with respect to $x$ will be:
\begin{equation}
    \left(\bar{\partial}_xI^\text{SR}\right)_{i,j} = \frac{1}{s^2} \sum_{k = 1}^{s^2} \left( I^\text{SR}_{i, j, k} - I^\text{SR}_{i, j + 1, k} \right).
    \label{e:mean_sr_derivative}
\end{equation}
It is very straightforward to generalize the equation~\ref{e:mean_sr_derivative} in all 8 directions.
Equation~\ref{e:mean_sr_derivative} is presented to illustrate the underlying principle, rather than to serve as a directly usable implementation.
In implementation, first, PixelUnshuffle is applied on $I^\text{SR}$ with a parameter scale factor $s$.
The latter produces $s^2$ channels, where in each channel, there are pixel neighbors with pixels that were originally in $I^\text{SR}$ in a distance of $s$.
Then, convolution with derivative kernels for each channel takes place.
Now, for each $s^2$ channel, we have up to 8 patches for derivatives in 8 possible directions.
Figure~\ref{fig:mean_sr_gradient} illustrates this process and provides some examples.
The last step is to calculate the mean across the $s^2$ channels, obtaining up to 8 mean SR derivatives.

\begin{landscape}
    \begin{figure}[]
        \centering
        \includegraphics[width=\linewidth]{images/mean_sr_gradient.pdf}
        \caption[Mean SR derivatives.]{Mean SR derivatives. Scale factor 3 was used as example. The derivative kernels are of size $3\times3$. In each kernel, all entries are zero except for the center pixel, which is set to 1, and the pixel in the corresponding directional position, which is set to -1. No padding is added in convolution with derivative kernels.}
        \label{fig:mean_sr_gradient}
    \end{figure}
\end{landscape}

\begin{landscape}
    \begin{table}[]
    \centering
    \begin{tabular}{|c|c|c|c|c|c|c|}
        \hline
         resampling method & conservation law & simple gradient & continuity & Sobel gradient & G-Loss & direction continuity \\ \hline \hline
         nearest        & 0.1343 & 0.3824 & 0.1635 & 0.6320 & 0.1907 & 0.0796 \\ \hline
         bilinear       & 0.0358 & 0.1024 & 0.0610 & 0.0646 & 0.0512 & 0.0048 \\ \hline
         cubic          & \textit{0.0225} & \textit{0.0609} & \textit{0.0016} & \textit{0.0199} & \textit{0.0303} & \textit{0.0029} \\ \hline
         cubic spline   & 0.0766 & 0.2038 & 0.1278 & 0.3754 & 0.1046 & 0.0171 \\ \hline
         Lanczos        & 0.0366 & 0.1000 & 0.0341 & 0.0864 & 0.0504 & 0.0064 \\ \hline
         average        & \textbf{0.0000} & \textbf{0.0000} & \textbf{0.0000} & \textbf{0.0000} & \textbf{0.0000} & \textbf{0.0000} \\ \hline
    \end{tabular}
    \caption[Resampling method experiment]{Resampling method experiment. Average violation of each constraint for different resampling methods between LR and HR images from ReKIS. Lowest, resp. second lowest violation for each constraint is highlighted as \textbf{bold}, resp. in \textit{italics}.}
    \label{tab:resampling_exp}
\end{table}
\end{landscape}


\begin{figure}[htbp]
    \centering
    \includegraphics[width=\linewidth]{images/sgl.pdf}
    \caption{Simple gradient loss.}
    \label{fig:sgl}
\end{figure}


In conclusion, we compiled a literature review; identified both soft and hard physical constraints --  specifically a conservation law -- and gathered five additional physically inspired loss-function augmentations, which we reformulated as soft constraints using the concept of dilated gradient consistency. 
We carried out experiments examining how the LR model input is obtained, focusing on different resampling methods. 
Furthermore, we investigated various soft constraint types, their weighting relative to the pixel-wise error, and the model sizes corresponding to unconstrained, soft-constrained, and hard-constrained settings.
We have demonstrated the positive effect of physical constraints on the model's performance.


In contrast, Structural Similarity (SSIM, Equation~\ref{e:ssim},~\cite{Wang2004}) primarily focuses on the structural characteristics of the image.
Furthermore, the Equation~\ref{e:ssim} is usually applied to images using a sliding window approach, with the default kernel size $11 \times 11$ in the PyTorch library.
The average of those partial values is then meant to be $\text{SSIM}{\left(Y,\hat{Y}\right)}$.
\begin{equation}
\text{SSIM}(x, y) \coloneqq \frac{(2\mu_x\mu_y + c_1)(2\sigma_{xy}+c_2)}{(\mu_x^2 + \mu_y^2+c_1)(\sigma_x^2+\sigma_y^2+c_2)},
\label{e:ssim}
\end{equation}
where $\mu_x, \mu_y$ are the pixel means of $x$ and $y$; $\sigma_x, \sigma_y, \sigma_{xy}$ are the variances of $x$ and $y$, respectively, as well as their covariance.
Moreover, for stability issues, when $(\mu_x^2 + \mu_y^2)$ or $(\sigma_x^2+\sigma_y^2)$ is close to zero, small constants $c_1, c_2$ are introduced.
Without loss of generality, their definition will be left out of this work and can be found in~\cite{Wang2004}.


$$
\text{PSNR}(Y, \hat{Y}) \coloneqq 10 \log_{10}\left(\frac{\text{MAX}^2}{\text{MSE}\left(Y, \hat{Y}\right)}\right)
$$

\noindent where $\text{MAX}$ is the maximum possible value of a pixel.
(TODO: Harder training.py 207: why is max target max? what if sr > target max??)



\begin{figure}[htbp]
    \centering
    \includegraphics[width=0.7\linewidth]{images/rekis_dev.pdf}
    \caption[ReKIS mean temperature deviation 1961--1991]{ReKIS mean temperature deviation hrough 1961--1991.}
    \label{fig:rekis_dev}
\end{figure}

Moreover, the differences between the means across all the spatiotemporal dimensions and the means through the temporal dimension (Figure~\ref{fig:rekis_mean}) are calculated and displayed in Figure~\ref{fig:rekis_dev}.
The latter indicates that the southern mountainous subregion has mean temperatures up to 6$^\circ$C lower than the overall regional mean, which is a much stronger deviation compared to most of the region, where the pixel mean is approximately 1--2$^\circ$C higher than the mean of the entire region.


Some existing studies fail to account for the fact that the variable being downscaled is a physical quantity and, as such, must obey the governing physical principles~\cite{Beucler2021}. 
The main objective, minimizing a standard pixel-wise metric (e.g., MSE, RMSE, MAE), becomes an obstacle for the model to learn physical relations, which is rather producing physically inconsistent outputs only to avoid penalty~\cite{Beucler2021}. 
Therefore, if we want to improve physical consistency, we need either to modify the loss function to penalize these physical errors (soft-constraining) or to adjust the model's architecture to completely forbid such behavior (hard-constraining).
Harder et al.~\cite{Harder2023} and Beucler et al.~\cite{Beucler2021} both provided a framework for implementing a constraint in the given form as soft- and hard-constraint. 
In both of these works, these techniques not only lead to physically (more) consistent outputs, but also the incorporation of appropriate physical constraints exerts a demonstrably positive effect on model accuracy.


In the work of Košťál et al.~\cite{Kostal2025}, using the same data, the best RMSE for the EDSR model on the validation set for a scale factor of 10 was 0.037.
For the best SwinIR model, this value was slightly lower at 0.032.
We have successfully improved these values to 0.0306 on test set while using a higher scale factor of 16.
Moreover, the best EDSR model in Košťál et al.~\cite{Kostal2025} consisted of 128 features and 64 residual blocks, thus having about 20 million parameters, similar to ours.
However, to reach the same performance, the hard-constrained model would need 20 times fewer parameters (for 64 features and 4 residual blocks; see Figure~\ref{fig:sizes}, Table~\ref{tab:sizes}).
